\section{An exact-charge relative theorem for screened lattice fermions}
\label{sl-sec}

This section establishes the screened lattice intermediate problem proposed
in Section~\ref{sec:reduction}. The result is a direct consequence of
bounded local occupation, a summable external potential, and a canonical
energy comparison proved below. It applies to the exact many-body energy
and to its exact-charge Slater restriction. No spectral-gap theorem or
small-coupling expansion is needed. The constants are uniform in volume,
but not in the removal of screening or the lattice ultraviolet regulator.
This is a benchmark theorem for a different, explicitly specified model;
it is not a theorem about continuum QED.

\subsection{Local Hamiltonians and canonical energy continuity}
\label{sl-local}

Let \(\Omega\) be a finite set of lattice sites, with \(2r\) fermion modes
per site, where \(r\) is a fixed positive integer. Put
\[
\mathcal F_\Omega=\bigwedge\bigl(\ell^2(\Omega)\otimes\mathbb C^{2r}\bigr),
\qquad D=2r|\Omega|,
\qquad n_x=\sum_{a=1}^{2r}c_{x,a}^*c_{x,a},
\qquad N=\sum_xn_x.
\]
Write \(N_* = r|\Omega|\) and \(Q=N-N_*\). The neutral vacuum energy will
always mean the minimum in the exact sector \(Q=0\), even if another
charge sector has a lower energy. Consider a decomposition
\begin{equation}
H_0=C_\Omega I+\sum_{\varnothing\ne X\subset\Omega}h_X,
\qquad h_X=h_X^*,\qquad [h_X,N_X]=0,
\qquad N_X=\sum_{x\in X}n_x,
\label{sl-local-h}
\end{equation}
where \(h_X\) belongs to the CAR algebra of the modes at sites in \(X\).
In particular every \(h_X\) is even. The scalar \(C_\Omega\) can include a
fixed free-vacuum normalization; its size is irrelevant below. Define
\begin{equation}
d_X=2r|X|,
\qquad
J=\sup_{x\in\Omega}\sum_{X\ni x}d_X\|h_X\|.
\label{sl-J}
\end{equation}
For a family of volumes the hypothesis is that \(J\) is bounded by a
volume-independent constant. Finite range is sufficient but not
necessary. Summable two-site interactions also satisfy this hypothesis.

For any charge-preserving Hamiltonian \(H\), let
\(E_{\rm MB}(n;H)\) be its minimum over all normalized vectors with
\(N=n\), and let \(E_{\rm SD}(n;H)\) be the minimum over normalized
Slater determinants in that sector. These minima exist: the Hilbert
space is finite dimensional and the relevant Grassmannian is compact.
The Slater class is an exact-particle-number class, not a mean-charge
relaxation.

\begin{lemma}[Uniform comparison between exact charge sectors]
\label{sl-charge-lip}
For \(\mathcal A\in\{{\rm MB},{\rm SD}\}\) and
\(0\le n,n'\le D\), the Hamiltonian \eqref{sl-local-h} satisfies
\begin{equation}
\left|E_{\mathcal A}(n';H_0)-E_{\mathcal A}(n;H_0)\right|
\le 2J|n'-n|.
\label{sl-charge-lip-eq}
\end{equation}
This assertion includes sectors arbitrarily close to the empty and fully
occupied sectors.
\end{lemma}

\begin{proof}
Use a single index \(i\) for a mode and, for each support \(X\), define
\[
A_X^+=\sum_{i\in X}c_i[h_X,c_i^*]
       =\sum_{i\in X}c_ih_Xc_i^*-(d_X-N_X)h_X,
\]
where the sum over \(i\in X\) includes all internal modes. Both terms are
self-adjoint and commute with \(N_X\). Moreover \(A_X^+\) vanishes on
the fully occupied local sector. Its norm is at most
\(2d_X\|h_X\|\). On every other local number sector,
\(d_X-N_X\ge1\), so the block decomposition by \(N_X\) proves the
operator inequality
\begin{equation}
-2d_X\|h_X\|(d_X-N_X)
\le A_X^+\le 2d_X\|h_X\|(d_X-N_X).
\label{sl-local-create}
\end{equation}
Similarly
\[
A_X^-=\sum_{i\in X}c_i^*[h_X,c_i]
       =\sum_{i\in X}c_i^*h_Xc_i-N_Xh_X
\]
is self-adjoint, commutes with \(N_X\), vanishes on the empty local
sector, and obeys
\begin{equation}
-2d_X\|h_X\|N_X\le A_X^-\le 2d_X\|h_X\|N_X.
\label{sl-local-annihilate}
\end{equation}
Evenness implies that modes outside \(X\) commute with \(h_X\). Thus
\eqref{sl-J}--\eqref{sl-local-annihilate} give
\begin{equation}
\sum_i c_i[H_0,c_i^*]\le 2J(D-N),
\qquad
\sum_i c_i^*[H_0,c_i]\le 2JN.
\label{sl-summed-comm}
\end{equation}
For example the first bound follows by writing
\(d_X-N_X=\sum_{x\in X}(2r-n_x)\) and collecting the coefficient
of each positive operator \(2r-n_x\).

Let \(\Psi\) minimize the energy in the \(n\)-particle class under
consideration. If \(n<D\), the creation trials satisfy
\[
\sum_i\|c_i^*\Psi\|^2=D-n,
\]
and, using \eqref{sl-summed-comm},
\begin{align*}
\sum_i\langle c_i^*\Psi,H_0c_i^*\Psi\rangle
&=(D-n)\langle\Psi,H_0\Psi\rangle
 +\left\langle\Psi,\sum_i c_i[H_0,c_i^*]\Psi\right\rangle\\
&\le(D-n)\bigl(E_{\mathcal A}(n;H_0)+2J\bigr).
\end{align*}
At least one trial is nonzero. Its energy, or equivalently the weighted
average of the nonzero trial energies, proves
\(E_{\mathcal A}(n+1;H_0)\le E_{\mathcal A}(n;H_0)+2J\).
The annihilation trials give
\(E_{\mathcal A}(n-1;H_0)\le E_{\mathcal A}(n;H_0)+2J\) for \(n>0\).
Applying these two inequalities to adjacent sectors and iterating proves
\eqref{sl-charge-lip-eq}.

For completeness, all these trials stay in the Slater class when
\(\Psi\) is a Slater determinant. If its occupied subspace is \(F\),
then \(c^*(f)\Psi\) is the wedge of \((1-P_F)f\) with \(\Psi\),
hence either zero or decomposable. For annihilation, rotate an
orthonormal basis of \(F\) so that its first vector is proportional
to \(P_Ff\); contraction by \(f\) then removes that first vector.
Thus \(c(f)\Psi\) is also zero or decomposable. The proof does not
require a Slater minimizer to be an eigenvector of \(H_0\).
\end{proof}

\subsection{Relative energy and relative correlation estimates}
\label{sl-relative}

Take integer nuclear charges \(1\le Z_i\le Z_{\max}\), and write
\(Z_{\rm tot}=\sum_iZ_i\). Suppose
\(0\le N_*+Z_{\rm tot}\le D\), the exact finite-mode admissibility
condition. Let
\begin{equation}
H_\nu=H_0+V_\nu+U_\nu I,
\qquad [V_\nu,N]=0,
\qquad \|V_\nu\|\le aZ_{\rm tot},
\qquad U_\nu\ge0,
\label{sl-source-assumption}
\end{equation}
with a constant \(a\) independent of volume, the number of sources,
and their positions. No disjointness of source supports is assumed.

\begin{theorem}[Screened lattice vacuum-relative stability]
\label{sl-relative-theorem}
Set \(C=2J+a\). For both \(\mathcal A={\rm MB}\) and
\(\mathcal A={\rm SD}\),
\begin{equation}
U_\nu-CZ_{\rm tot}
\le E_{\mathcal A}(N_*+Z_{\rm tot};H_\nu)
      -E_{\mathcal A}(N_*;H_0)
\le U_\nu+CZ_{\rm tot}.
\label{sl-two-sided}
\end{equation}
In particular, with the actual interacting neutral vacuum
\(E_0^{\rm vac}=E_{\rm MB}(N_*;H_0)\),
\begin{equation}
E_{\rm MB}(N_*+Z_{\rm tot};H_\nu)-E_0^{\rm vac}
\ge-CZ_{\max}K.
\label{sl-stability}
\end{equation}
Define the nonnegative Slater correlation errors
\begin{align*}
c_\nu&=E_{\rm SD}(N_*+Z_{\rm tot};H_\nu)
       -E_{\rm MB}(N_*+Z_{\rm tot};H_\nu),\\
c_0&=E_{\rm SD}(N_*;H_0)-E_{\rm MB}(N_*;H_0).
\end{align*}
Then the following two-sided relative correlation estimate also holds:
\begin{equation}
|c_\nu-c_0|\le 2CZ_{\rm tot}\le 2CZ_{\max}K.
\label{sl-relative-correlation}
\end{equation}
\end{theorem}

\begin{proof}
For any fixed particle number, the operator inequalities
\(-aZ_{\rm tot}I\le V_\nu\le aZ_{\rm tot}I\) bound the change of
either variational infimum by \(aZ_{\rm tot}\). Apply
Lemma~\ref{sl-charge-lip} to compare the remaining source-free energy
in sector \(N_*+Z_{\rm tot}\) with that in sector \(N_*\). This proves
\eqref{sl-two-sided}; dropping \(U_\nu\ge0\) gives
\eqref{sl-stability}. Subtract the two relative energy differences
in \eqref{sl-two-sided}. The same specified nuclear repulsion occurs
in both and cancels, leaving \eqref{sl-relative-correlation}.
No state-dependent or source-dependent vacuum subtraction is used.
\end{proof}

The theorem also covers bounded local defects:
\(V_\nu=\sum_iB_i\), with \([B_i,N]=0\) and
\(\|B_i\|\le aZ_i\), regardless of overlap between their supports.
The charge-sector lemma is essential if the physical source sector
differs from the neutral vacuum sector. Simply comparing both
Hamiltonians with their unconstrained ground energies would not prove
\eqref{sl-stability} with the stated vacuum.

\subsection{A concrete Yukawa model with particle--hole pair terms}
\label{sl-yukawa}

Fix lattice spacing one, \(\ell\ge3\), and
\(\Omega_\ell=(\mathbb Z/\ell\mathbb Z)^3\). Its Laplacian is
\[
(\Delta_\ell f)(x)=\sum_{j=1}^3
\bigl(f(x+e_j)+f(x-e_j)-2f(x)\bigr).
\]
For \(\mu>0\), define the translation-invariant screened kernel by
\begin{equation}
W_{\mu,\ell}(x-y)
=4\pi\langle\delta_x,(-\Delta_\ell+\mu^2)^{-1}\delta_y\rangle.
\label{sl-yukawa-kernel}
\end{equation}
The zero Fourier mode is retained. This is a screened Poisson
resolvent, not the zero-mode-deleted Coulomb Green function. The finite
matrix identity
\[
(-\Delta_\ell+\mu^2)^{-1}
=\int_0^\infty e^{-\mu^2t}e^{t\Delta_\ell}\,dt
\]
proves positivity of every kernel entry: the heat semigroup has
nonnegative entries, as follows by expanding the exponential of the
nonnegative adjacency matrix. Since \(e^{t\Delta_\ell}1=1\), it also
proves the exact, volume-independent row sum
\begin{equation}
W_{\mu,\ell}(x)\ge0,
\qquad
S_\mu:=\sum_{x\in\Omega_\ell}W_{\mu,\ell}(x)
=\frac{4\pi}{\mu^2}.
\label{sl-row-mass}
\end{equation}

Put \(q_x=n_x-r\), so \(\|q_x\|=r\), and let
\begin{equation}
H_0=H_{\rm kin}
 +\frac g2\sum_{x,y\in\Omega_\ell}
 W_{\mu,\ell}(x-y)q_xq_y,
\qquad g\ge0.
\label{sl-concrete-h0}
\end{equation}
The matter interaction includes its specified finite on-site term.
All operators are literal finite CAR operators; no additional Wick
prescription is implicit. Suppose \(H_{\rm kin}\) has a decomposition
of the form \eqref{sl-local-h} with constant \(J_{\rm kin}\).
For the off-site term indexed by \(\{x,y\}\),
\(\|gW(x-y)q_xq_y\|\le gr^2W(x-y)\) and \(d_{\{x,y\}}=4r\).
For the on-site term the corresponding bound is
\(gr^2W(0)/2\) and \(d_{\{x\}}=2r\). Therefore
\begin{equation}
J\le J_{\rm kin}+4gr^3S_\mu.
\label{sl-J-concrete}
\end{equation}

One concrete kinetic choice, with \(r=2\), is the lattice Dirac band
Hamiltonian
\begin{equation}
h_W(p)=\sum_{j=1}^3\alpha_j\sin p_j
 +\beta\left(m+w\sum_{j=1}^3(1-\cos p_j)\right),
\qquad m>0,\quad w\ge0,
\label{sl-wilson}
\end{equation}
and
\[
H_{\rm kin}=d\Gamma(h_W)-\operatorname{Tr}(h_WP_-),
\qquad P_-=\boldsymbol1_{(-\infty,0)}(h_W).
\]
Its bands are \(\pm\varepsilon(p)\), each with multiplicity two,
and \(\varepsilon(p)\ge m\). Filling the negative bands gives
\(N_*=2|\Omega_\ell|\). This kinetic Hamiltonian is nonnegative,
and the scalar subtraction does not enter \(J_{\rm kin}\).
For an entirely explicit, deliberately loose bound, use the on-site
term \(c_x^*\beta(m+3w)c_x\) and nearest-neighbor matrices
\(T_j=-i\alpha_j/2-w\beta/2\). The on-site norm is
\(2(m+3w)\). A bond acts on eight modes and has norm at most
\(8\|T_j\|\le4(1+w)\); at most six bonds meet each site.
Consequently
\begin{equation}
J_{\rm kin}\le8(m+3w)+192(1+w).
\label{sl-J-kinetic}
\end{equation}

In the particle--hole representation of the two bands, positive-band
annihilators are electron annihilators \(b\), and negative-band
annihilators are positron creators \(d^*\). Then
\(Q=N-N_*=N_e-N_p\), while \(H_{\rm kin}\) is the positive free
electron--positron Hamiltonian. Local densities and their interactions
generically contain \(b^*d^*\) and \(db\) terms in this representation.
These terms are retained; neither electron number nor positron number
is fixed. Charge conservation fixes only their difference. The theorem
does not rely on conservation of the number of pairs.

Choose sites \(R_i\in\Omega_\ell\), with repetitions permitted, and
the lattice point distribution
\(\nu(x)=\sum_iZ_i\delta_{R_i,x}\). Define
\begin{align}
V_\nu&=-g\sum_iZ_i\sum_xW_{\mu,\ell}(x-R_i)q_x,
\label{sl-source-potential}\\
U_\nu&=g\sum_{i<j}Z_iZ_jW_{\mu,\ell}(R_i-R_j)\ge0.
\label{sl-nuclear-repulsion}
\end{align}
Each individual nuclear self-energy
\(gZ_i^2W_{\mu,\ell}(0)/2\) is excluded by this explicit convention.
It is finite on the lattice. Distinct sources at the same site still
contribute their mutual repulsion in \eqref{sl-nuclear-repulsion}.
These are discrete point sources; the finite on-site potential does
not represent the singularity of a continuum point nucleus.

Equations~\eqref{sl-row-mass} and \eqref{sl-source-potential} give
\(\|V_\nu\|\le grS_\mu Z_{\rm tot}\). Thus a valid constant in
Theorem~\ref{sl-relative-theorem} is
\begin{equation}
C_\mu=2J_{\rm kin}+g(8r^3+r)\frac{4\pi}{\mu^2}.
\label{sl-Cmu}
\end{equation}
In particular
\[
E_{\rm MB}(N_*+Z_{\rm tot};H_\nu)-E_0^{\rm vac}
\ge-C_\mu Z_{\max}K,
\qquad
|c_\nu-c_0|\le2C_\mu Z_{\max}K.
\]
These bounds hold at arbitrary \(g\ge0\), for every admissible exact
charge sector and every source configuration, uniformly in \(\ell\).
There are no photons or local Gauss-law constraints in this purely
fermionic model. Exact neutrality is the specified physical charge
sector, rather than a conclusion of a gauge constraint.

\subsection{What the theorem establishes, and the Coulomb obstruction}
\label{sl-scope}

The result is stronger than finite-dimensional boundedness: the
comparison is with the interacting neutral vacuum and its constant is
independent of volume, source overlap, and pair number. It also proves
the screened analogue of the sufficient relative correlation estimate
\eqref{eq:rel-corr}. Its mechanism is nevertheless elementary: the
entire external perturbation has norm of order \(K\). It supplies no
decay estimate for the response of the vacuum and no refined expansion
of the correlation energy. The bound need not vanish when the
interaction is removed, although the actual correlation error then
vanishes for a quadratic Hamiltonian.

No source-free gap, nor persistence of a gap after arbitrary sources
are added, enters the proof. The gap results of De Roeck--Salmhofer
\cite{DRS2019} can be investigated separately for suitable weakly
interacting lattice systems. They are not needed for this relative
energy estimate, and do not by themselves provide uniform gaps or
localized response in arbitrary nuclear backgrounds.

The coefficient obstruction to recovering an expanding Coulomb window
is explicit. Here \(S_\mu=4\pi/\mu^2\), so the estimate deteriorates
as the screening length grows. For the unscreened mean-zero discrete
Coulomb kernel \(G_\ell\), its Fourier coefficient at
\(k=(2\pi/\ell,0,0)\) gives
\begin{equation}
\sum_x|G_\ell(x)|
\ge\left|\sum_xG_\ell(x)e^{-ik\cdot x}\right|
=\frac{4\pi}{4\sin^2(\pi/\ell)},
\label{sl-coulomb-l1}
\end{equation}
which grows quadratically in the box side. Likewise, the absolute
sum of a Coulomb kernel retained only out to distance \(R\) grows
as \(R^2\) in three dimensions. This absolute-row-sum estimate
therefore cannot yield a uniform coefficient while the
Coulomb window expands. Global neutrality alone does not alter these
absolute-norm estimates. A new cancellation or electrostatic estimate
would be required; \eqref{sl-coulomb-l1} is an obstruction to this
proof method, not a counterexample to Coulomb stability.

Finally the lattice spacing has been fixed throughout. Restoring a
variable spacing changes the kinetic norms, the source potential, and
the density of available modes. No uniform ultraviolet limit, charge
renormalization, continuum point-source theorem, or quantized-photon
stability follows from \eqref{sl-Cmu}. The successful screened theorem
removes the need for a gap-based argument at this preliminary stage;
it identifies the nonsummable Coulomb tail and continuum ultraviolet
limit as the places where stronger estimates are needed.
